%
% Author: Dylan Gormley | dylan.gormley@mail.wvu.edu
% This is beamer's amurmaple theme configured with WVU colors and logos.
%
% 1) Must be compiled with with LuaTex
% 2) Must be set to aspectratio=169
%
\documentclass[aspectratio=169,10pt]{beamer}

%
% Adjust theme options, choose secondary colors and set up title page in styles/amurmaplewvu.sty
%
\usepackage{amurmaplewvu}

%
% The example below is taken from the theme author's repository, with a few changes
% https://gitlab.gutenberg-asso.fr/mchupin/amurmaple/-/blob/main/test/beamer-amurmaple-black.tex
%
%\setbeameroption{show notes on second screen }
\usepackage{pgfpages}
\usepackage{lipsum}

\pdfstringdefDisableCommands{%
  \DeclareRobustCommand{\translate}[1]{#1}%
}


\begin{document}
\maketitle

\sepframe[title={Table of contents}]

\frame{\tableofcontents}

\section{The first section}

\subsection{See the documentation of the beamer class for details}

%%%%%%%%%%%%%%%%%%%%%%%%%%%%%%%%
% ----------- FRAME ------------
%%%%%%%%%%%%%%%%%%%%%%%%%%%%%%%%
\begin{frame}[<+->]{Lorem ipsum}{Sed commodo posuere pede}
    \begin{block}{Block title}
    Cras viverra metus rhoncus sem.
    \end{block}

    \begin{itemize}
    	\item Lorem ipsum dolor sit amet
    	\item Consectetuer adipiscing elit
    	\item Ut purus elit, vestibulum ut
    \end{itemize}

\begin{quotation}[Donald E. Knuth, \emph{The \TeX book}]
Gentle reader: This is a handbook about \TeX, a new typesetting system G intended for the creation of beautiful books—and especially for books that contain a lot of mathematics.
\end{quotation}

\begin{quotation}
  Gentle reader: This is a handbook about \TeX, a new typesetting system G intended for the creation of beautiful books—and especially for books that contain a lot of mathematics.
  \end{quotation}
  
\end{frame}

\section{Second section}

\sepframe

\begin{frame}{Test}
  \framesubtitle{Subtitle of the frame}
  Hello
  \begin{block}{Answered Questions}
    How many primes are there?
\end{block}
\begin{block}{Open Questions}
Is every even number the sum of two primes?
\end{block}
\begin{remark}[complements]
Is this a relevant remark?
\end{remark}
\end{frame}

\begin{frame}
  \framesection{Paragraph}
  \beamerbutton{Button}\beamerskipbutton{Skip Button}\beamerreturnbutton{Return}
  \par
  \structure{Test structure}
  \alert{Test alert}
  \boxalert{Test boxalert}

  \begin{abstract}
    This is an abstract.
  \end{abstract}

  \begin{information}
    This is an information.
  \end{information}
\end{frame}

\begin{frame}
  \begin{block}{Block}
    Test block
  \end{block}
  \begin{alertblock}{Alert Block}
    Test block
  \end{alertblock}
  \begin{exampleblock}{Example Block}
    Test block
  \end{exampleblock}
\end{frame}



\begin{frame}
\frametitle{A Theorem on Infinite Sets}
\begin{theorem}<1->
There exists an infinite set.
\end{theorem}
\begin{proof}<2->
This follows from the axiom of infinity.
\end{proof}
\begin{example}<3->[Natural Numbers]
The set of natural numbers is infinite.
\end{example}
\begin{definition}[Title of def.]
Test
\end{definition}
\begin{corollary}
Test
\end{corollary}
\end{frame}

\begin{frame}
  \begin{beamercolorbox}{beamer color}
Text\footnote{On a fast machine.}
\end{beamercolorbox}
\end{frame}

\section{New one}

\begin{frame}
\sectionpage
\end{frame}


\part{Name of the part}
\begin{frame}
\partpage
\end{frame}



\begin{frame}
\begin{itemize}
\item<1-> Eggs
\item<2-| alert@2> Plants
  \begin{itemize}
  \item test
  \end{itemize}
\note[item]<2>{Tell joke about plants.}
\note[item]<2>{Make it short.}
\item<3-> Animals
\end{itemize}
\begin{enumerate}
\item<1-> Eggs
\item<2-| alert@2> Plants
\note[item]<2>{Tell joke about plants.}
\note[item]<2>{Make it short.}
\item<3-> Animals
  \begin{enumerate}
  \item Dog
  \item Cat
  \end{enumerate}
\end{enumerate}
\end{frame}

\begin{frame}{2 columns}
  \begin{columns}[t]
    \begin{column}{5cm}
      Two\\lines.
    \end{column}
    \begin{column}{5cm}
      One line (but aligned).
    \end{column}
  \end{columns}
\end{frame}



\thanksframe{Thank you for your attention!}

\appendix
\section{\appendixname}
\section{Appendix Name}

\begin{frame}
\tableofcontents
\end{frame}

\subsection{Additional material}

\begin{frame}
Details
\end{frame}


\end{document}
